% Preamble requirements:
% \usepackage{amsmath,amssymb,booktabs,array,enumitem,tikz}
% \usetikzlibrary{arrows.meta,positioning,shapes.geometric}

\section{Problem Formulation and System Overview}
\label{sec:problem_formulation}

This thesis formulates flow-level network intrusion detection as a hierarchical
binary classification problem. The hierarchy arises from two nested structures
in network traffic: packets form bidirectional flows, and related flows form
larger temporal or host-level behaviors. Conventional flow-based classifiers
typically compress the first structure into aggregate statistics and classify
each resulting record independently. This is efficient, but it can discard
packet order and irregular timing and can overlook evidence distributed across
several flows \cite{SommerPaxson2010,Han2023GTID,LiuPatras2022NetSentry}.
Transformer-based NIDS research has demonstrated the value of both flow
sequences and long-range dependencies
\cite{Manocchio2024FlowTransformer,MartinReizabal2026FlowWindow}; however, the
two levels of the hierarchy are not always modeled jointly. The objective of
this work is therefore to learn a representation that first captures
\emph{intra-flow} packet dynamics and then incorporates \emph{inter-flow}
context for the final decision.

\subsection{Detection Setting and Design Requirements}
\label{subsec:detection_setting}

The detector operates on traffic extracted from PCAP traces. A custom Suricata
output plugin produces packet-level and flow-level records that can be aligned
through a flow identifier. Suricata is used as a practical traffic-processing
and logging engine \cite{Roesch1999,SuricataGuide2026}; its alert outcome is
used only during label construction. Alert counts, signature identifiers, rule
messages, and other direct indicators of the label are excluded from the model
input. The proposed detector also avoids direct payload content and instead
uses packet-header fields, packet direction and timing, aggregate flow
statistics, and metadata required to construct flow relationships. This setting
supports encrypted traffic because it does not require payload decryption.

The formulation is guided by the following requirements:

\begin{enumerate}[label=\textbf{R\arabic*:},leftmargin=*]
    \item \textbf{Hierarchical representation.} The model should preserve the
    packet sequence inside each flow and should also represent dependencies
    among related flows.

    \item \textbf{Explicit network time.} Packet position and elapsed time must
    be distinguishable. Adjacent packets can be separated by very different
    intervals, so ordinal position alone is insufficient
    \cite{Han2023GTID}.

    \item \textbf{Payload independence.} Predictions should not require raw
    application payloads, thereby reducing dependence on plaintext visibility
    and limiting exposure of content-sensitive data.

    \item \textbf{Causal context.} The context of a target flow may contain the
    target itself and eligible historical flows, but it must not contain future
    flows or labels. This requirement makes the formulation compatible with an
    online detector.

    \item \textbf{Imbalance-aware learning and evaluation.} Because malicious
    flows are the minority class, the learning objective and reported metrics
    must not be dominated by benign traffic.

    \item \textbf{Reproducible separation of data.} Preprocessing, context
    construction, model selection, and threshold selection must respect the
    training, validation, and test partitions. This limits common sources of
    experimental leakage in security machine learning
    \cite{Arp2022DosDonts,Ring2019Datasets}.
\end{enumerate}

The scope of the thesis is supervised binary flow classification rather than
attack prevention or rule generation. The output is a probability that the
target flow is malicious. Multiclass attack attribution, zero-day guarantees,
and automated response are outside the formal scope and are discussed as future
work.

\subsection{System Overview}
\label{subsec:system_overview}

Figure~\ref{fig:system_overview} summarizes the complete processing path. The
pipeline begins with raw PCAP traffic and ends with one prediction per target
flow. Packet and flow records are extracted together, aligned, cleaned, and
partitioned before model training. Stage1 transforms each packet sequence into
a time-aware representation and combines it with the corresponding statistical
flow representation through late fusion. The resulting Stage1 embedding is a
single fixed-dimensional vector per flow. Stage2 constructs a causal sequence
of these embeddings according to a selected context relation and predicts the
label of the target flow.

\begin{figure}[t]
\centering
\begin{tikzpicture}[
    node distance=4.5mm,
    block/.style={
        rectangle,
        draw,
        rounded corners=2pt,
        align=center,
        text width=11.3cm,
        minimum height=7.5mm,
        fill=black!3
    },
    arrow/.style={-{Latex[length=2mm]},thick}
]
    \node[block] (pcap) {Raw PCAP traffic};
    \node[block, below=of pcap] (extract) {Custom Suricata output plugin};
    \node[block, below=of extract] (records) {Packet records, flow statistics, flow identifiers, timestamps, and endpoint metadata};
    \node[block, below=of records] (prep) {Alignment, cleaning, chronological ordering, feature transformation, and train/validation/test partitioning};
    \node[block, below=of prep] (stage1) {Stage1: time-aware packet-sequence Transformer, attention pooling, and late fusion with encoded flow statistics};
    \node[block, below=of stage1] (embedding) {One fused Stage1 embedding per flow};
    \node[block, below=of embedding] (context) {Causal context construction: time-only, source-host, destination-host, or endpoint relation};
    \node[block, below=of context] (stage2) {Stage2: context encoder and target-flow classifier};
    \node[block, below=of stage2] (output) {Malicious-flow probability and binary prediction};

    \draw[arrow] (pcap) -- (extract);
    \draw[arrow] (extract) -- (records);
    \draw[arrow] (records) -- (prep);
    \draw[arrow] (prep) -- (stage1);
    \draw[arrow] (stage1) -- (embedding);
    \draw[arrow] (embedding) -- (context);
    \draw[arrow] (context) -- (stage2);
    \draw[arrow] (stage2) -- (output);
\end{tikzpicture}
\caption{Overview of the proposed hierarchical two-stage intrusion-detection
pipeline. Labels and Suricata alert metadata are not used as model features.}
\label{fig:system_overview}
\end{figure}

The two stages are trained sequentially in the current system. Stage1 is first
optimized for flow-level classification and then used to export fused
embeddings. Stage2 is trained on the exported embeddings and does not
back-propagate gradients into Stage1. This separation permits independent
evaluation of the representation learned in Stage1 and the contextual gain
introduced by Stage2.

\subsection{Input Space and Notation}
\label{subsec:input_notation}

Let the dataset contain $N$ bidirectional flows sorted by flow start time. The
$i$th flow is represented by

\begin{equation}
\mathcal{F}_i =
\left(
    \mathbf{X}_i,
    \boldsymbol{\tau}_i,
    \mathbf{b}_i,
    \mathbf{s}_i,
    \mathbf{q}_i,
    y_i
\right),
\end{equation}

where the components are defined in
Table~\ref{tab:problem_notation}. The complete supervised dataset is

\begin{equation}
\mathcal{D}=\{\mathcal{F}_i\}_{i=1}^{N}.
\end{equation}

\begin{table}[t]
\centering
\caption{Notation used in the hierarchical problem formulation.}
\label{tab:problem_notation}
\small
\setlength{\tabcolsep}{5pt}
\renewcommand{\arraystretch}{1.15}
\begin{tabular}{@{}p{2.3cm}p{3.2cm}p{9.7cm}@{}}
\toprule
\textbf{Symbol} & \textbf{Domain} & \textbf{Meaning} \\
\midrule
$N$ & $\mathbb{N}$ & Number of flows. \\
$T_i$ & $\mathbb{N}$ & Number of observed packets in flow $i$ before sequence selection. \\
$L_{\max}$ & $\mathbb{N}$ & Maximum number of packet tokens retained for Stage1. \\
$\mathbf{X}_i$ & $\mathbb{R}^{L_i\times d_p}$ & Selected packet-feature sequence, where $L_i\leq L_{\max}$ and $d_p$ is the packet-feature dimension. \\
$\boldsymbol{\tau}_i$ & $\mathbb{R}^{L_i}$ & Packet timestamps or elapsed-time values aligned with $\mathbf{X}_i$. \\
$\mathbf{b}_i$ & $\{0,1\}^{L_{\max}}$ & Packet mask; one denotes an observed packet and zero denotes padding. \\
$\mathbf{s}_i$ & $\mathbb{R}^{d_f}$ & Flow-level statistical feature vector. \\
$\mathbf{q}_i$ & $\mathcal{Q}$ & Context metadata: source identifier, destination identifier, flow start time, and split membership. \\
$y_i$ & $\{0,1\}$ & Flow label, with zero for benign and one for malicious. \\
$\mathbf{z}^{(1)}_i$ & $\mathbb{R}^{d_1}$ & Fused Stage1 embedding for flow $i$. \\
$\mathcal{C}^{(r)}_i$ & Subset of $\{1,\ldots,i-1\}$ & Historical context indices for target flow $i$ under relation $r$. \\
$W$ & $\mathbb{N}$ & Maximum Stage2 sequence length, including the target flow when it is appended. \\
\bottomrule
\end{tabular}
\end{table}

Before padding, flow $i$ contains an ordered packet sequence

\begin{equation}
\mathcal{P}_i=
\left(p_{i,1},p_{i,2},\ldots,p_{i,T_i}\right),
\qquad
\tau_{i,1}\leq\tau_{i,2}\leq\cdots\leq\tau_{i,T_i}.
\end{equation}

Each retained packet is represented by a feature vector
$\mathbf{x}_{i,t}\in\mathbb{R}^{d_p}$. Packet timing can be expressed through
both elapsed time from the first retained packet and local inter-arrival time:

\begin{align}
e_{i,t} &= \tau_{i,t}-\tau_{i,1}, \\
\Delta\tau_{i,t} &=
\begin{cases}
0, & t=1,\\
\max(0,\tau_{i,t}-\tau_{i,t-1}), & t>1.
\end{cases}
\end{align}

A deterministic sequence-selection operator $\Pi_{L_{\max}}$ converts a
variable-length flow into at most $L_{\max}$ packet tokens:

\begin{equation}
\mathbf{X}_i=
\Pi_{L_{\max}}
\left(
    \mathbf{x}_{i,1},\ldots,\mathbf{x}_{i,T_i}
\right).
\end{equation}

The operator covers the evaluated head, head--tail, and random strategies. Its
choice changes which packets are retained but does not change the label or the
flow-level statistical vector. Padding is applied only after selection and is
excluded from attention and pooling through $\mathbf{b}_i$.

The context metadata can be written as

\begin{equation}
\mathbf{q}_i=(u_i,v_i,t_i^{0},\rho_i),
\end{equation}

where $u_i$ and $v_i$ are the source and destination identifiers, $t_i^{0}$ is
the flow start time, and $\rho_i\in\{\mathrm{train},\mathrm{val},\mathrm{test}\}$
denotes split membership. These identifiers determine context membership; they
are not labels and need not be directly embedded as numerical packet features.

\subsection{Stage1: Intra-Flow Representation Learning}
\label{subsec:stage1_formulation}

Stage1 learns one representation from the packets of a single flow. Packet
features are first projected into a $d_1$-dimensional latent space:

\begin{equation}
\mathbf{e}_{i,t}=W_p\mathbf{x}_{i,t}+\mathbf{c}_p.
\end{equation}

The model then adds information about ordinal position and network time:

\begin{equation}
\widetilde{\mathbf{e}}_{i,t}
=
\mathbf{e}_{i,t}
+\phi_{\mathrm{pos}}(t)
+\phi_{\mathrm{time}}(e_{i,t},\Delta\tau_{i,t}).
\end{equation}

The functions $\phi_{\mathrm{pos}}$ and $\phi_{\mathrm{time}}$ are written
separately because packet index and elapsed time convey different information.
The exact scaling, transformation, and fusion of the two encodings are specified
in the methodology chapter. A mask-aware Transformer encoder
\cite{Vaswani2017} produces contextualized packet states:

\begin{equation}
\mathbf{H}^{p}_i
=
\mathrm{Transformer}_{1}
\left(
    \widetilde{\mathbf{E}}_i;
    \mathbf{b}_i
\right),
\qquad
\mathbf{H}^{p}_i\in\mathbb{R}^{L_{\max}\times d_1}.
\end{equation}

The valid packet states are reduced to one packet-sequence representation by
attention pooling. For each valid position,

\begin{align}
r_{i,t} &=
\mathbf{w}_{a}^{\top}
\tanh(W_a\mathbf{h}^{p}_{i,t}+\mathbf{c}_a),\\
a_{i,t} &=
\frac{
    b_{i,t}\exp(r_{i,t})
}{
    \sum_{k=1}^{L_{\max}}b_{i,k}\exp(r_{i,k})
},\\
\mathbf{z}^{p}_i &=
\sum_{t=1}^{L_{\max}}a_{i,t}\mathbf{h}^{p}_{i,t}.
\end{align}

In parallel, the statistical flow vector is mapped into the same latent space:

\begin{equation}
\mathbf{z}^{f}_i=E_f(\mathbf{s}_i).
\end{equation}

Late fusion produces the final Stage1 embedding:

\begin{equation}
\mathbf{z}^{(1)}_i
=
\mathrm{Fuse}_{\psi}
\left(
    \mathbf{z}^{p}_i,
    \mathbf{z}^{f}_i
\right).
\label{eq:stage1_fusion}
\end{equation}

This formulation is deliberately different from repeating
$\mathbf{s}_i$ at every packet position. Packet-sequence learning is completed
before the aggregate statistics are fused, allowing the contribution of each
view to be measured independently. The Stage1 classifier is

\begin{equation}
\mathbf{p}^{(1)}_i
=
\mathrm{softmax}
\left(
    h_1(\mathbf{z}^{(1)}_i)
\right),
\end{equation}

where $p^{(1)}_{i,1}$ is the predicted probability of a malicious flow. This
prediction supports direct evaluation of Stage1 and therefore operationalizes
RQ1 before any cross-flow context is introduced.

\subsection{Stage2: Causal Inter-Flow Context Modeling}
\label{subsec:stage2_formulation}

After Stage1 inference, every flow has an embedding
$\mathbf{z}^{(1)}_i$. Let the complete set of strictly historical flows for
target $i$ be

\begin{equation}
\mathcal{H}_i=\{j\mid j<i\},
\end{equation}

where indices follow a stable chronological ordering. The system defines a
binary relation $R_r(i,j)$ for each context method $r$:

\begin{align}
R_{\mathrm{time}}(i,j)
&=1,\\
R_{\mathrm{src}}(i,j)
&=\mathbb{I}[u_i=u_j],\\
R_{\mathrm{dst}}(i,j)
&=\mathbb{I}[v_i=v_j],\\
R_{\mathrm{endpoint}}(i,j)
&=\mathbb{I}
\left[
    \{u_i,v_i\}\cap\{u_j,v_j\}\neq\varnothing
\right].
\end{align}

The eligible historical set is

\begin{equation}
\mathcal{A}^{(r)}_i
=
\left\{
    j\in\mathcal{H}_i
    \mid
    R_r(i,j)=1
\right\}.
\end{equation}

Given a maximum Stage2 length $W$, the context builder retains at most the
$W-1$ most recent eligible and deduplicated flows:

\begin{equation}
\mathcal{C}^{(r)}_i
=
\mathrm{Recent}_{W-1}
\left(
    \mathcal{A}^{(r)}_i
\right).
\label{eq:context_set}
\end{equation}

If
$\mathcal{C}^{(r)}_i=(j_1,\ldots,j_{K_i})$ with $K_i\leq W-1$, the Stage2 input
sequence is

\begin{equation}
\mathbf{Z}^{(r)}_i
=
\left[
    \mathbf{z}^{(1)}_{j_1},
    \ldots,
    \mathbf{z}^{(1)}_{j_{K_i}},
    \mathbf{z}^{(1)}_i
\right].
\label{eq:stage2_sequence}
\end{equation}

The current flow is appended as the final valid token. It does not reveal its
label; it provides the representation that Stage2 must classify. Padding and a
context mask are added when fewer than $W-1$ historical flows are available.
Because the histories are updated only after the current context is built,
Equations~\ref{eq:context_set}--\ref{eq:stage2_sequence} exclude both the future
and accidental self-duplication.

The primary formulation uses an online context policy: earlier flows can remain
available as state when the detector moves from training to validation and test
periods, but no label is retrieved and no future flow is visible. A
split-isolated policy can additionally restrict candidates to the target split
for sensitivity analysis. The policy used in each experiment must be reported
because context availability and the risk of leakage depend on it.

Stage2 projects and positionally encodes the context sequence before applying a
second sequence encoder:

\begin{equation}
\mathbf{G}^{(r)}_i
=
\mathrm{Transformer}_{2}
\left(
    \mathrm{Proj}_{2}(\mathbf{Z}^{(r)}_i)
    +\Phi_{2};
    \boldsymbol{\mu}_i
\right),
\end{equation}

where $\boldsymbol{\mu}_i$ is the context mask. Let
$\mathbf{g}^{(r)}_i$ denote the state at the final valid, target-flow position.
The context-aware probability is

\begin{equation}
\mathbf{p}^{(2,r)}_i
=
\mathrm{softmax}
\left(
    h_2(\mathbf{g}^{(r)}_i)
\right).
\label{eq:stage2_prediction}
\end{equation}

A no-context classifier provides the control condition

\begin{equation}
\mathbf{p}^{(\mathrm{nc})}_i
=
\mathrm{softmax}
\left(
    h_{\mathrm{nc}}(\mathbf{z}^{(1)}_i)
\right).
\end{equation}

Comparing Equation~\ref{eq:stage2_prediction} with the no-context control tests
whether neighboring flows add information beyond the capacity of another
classifier. Comparing the four relations tests which temporal or endpoint
assumption is most useful, thereby operationalizing RQ2.

\subsection{Learning Objectives and Decision Rule}
\label{subsec:learning_objectives}

Both stages solve an imbalanced binary classification problem. For stage
$\ell\in\{1,2\}$, let $p^{(\ell)}_{i,y_i}$ denote the probability assigned to
the correct class. The focal objective is

\begin{equation}
\mathcal{L}_{\ell}
=
-\frac{1}{|\mathcal{D}_{\mathrm{train}}|}
\sum_{i\in\mathcal{D}_{\mathrm{train}}}
\alpha^{(\ell)}_{y_i}
\left(
    1-p^{(\ell)}_{i,y_i}
\right)^{\gamma_{\ell}}
\log p^{(\ell)}_{i,y_i},
\label{eq:focal_objective}
\end{equation}

where $\gamma_{\ell}\geq0$ controls the emphasis on difficult examples and
$\alpha^{(\ell)}_{y_i}$ optionally controls class weighting
\cite{Lin2017FocalLoss}. Cross-entropy is recovered when
$\gamma_{\ell}=0$ and the class weights are uniform. The precise loss
configuration is treated as an experimental setting rather than an inherent
requirement of the architecture.

The current implementation uses sequential optimization. First,

\begin{equation}
\theta_1^{*}
=
\arg\min_{\theta_1}
\mathcal{L}_{1}(\theta_1),
\end{equation}

after which Stage1 exports
$\{\mathbf{z}^{(1)}_i(\theta_1^{*})\}_{i=1}^{N}$. Stage2 then solves

\begin{equation}
\theta_2^{*}
=
\arg\min_{\theta_2}
\mathcal{L}_{2}
\left(
    \theta_2;
    \{\mathbf{z}^{(1)}_i(\theta_1^{*})\}_{i=1}^{N}
\right).
\end{equation}

This formulation prevents Stage2 experiments from silently changing the
Stage1 representation and permits controlled comparisons among context methods
and Stage2 architectures.

For the malicious class, the final probability is
$p_i=p^{(2,r)}_{i,1}$, or $p_i=p^{(1)}_{i,1}$ when Stage1 is evaluated alone. A
binary decision is obtained using

\begin{equation}
\widehat{y}_i
=
\mathbb{I}[p_i\geq\delta^{*}],
\end{equation}

where the decision threshold is selected exclusively on the validation set:

\begin{equation}
\delta^{*}
=
\arg\max_{\delta\in\Delta}
M_{\mathrm{val}}(\delta).
\end{equation}

Here, $M_{\mathrm{val}}$ is the predefined threshold-selection metric and
$\Delta$ is the threshold search grid. The test set is evaluated once with the
selected model and threshold; it is not used to tune either quantity.

\subsection{Mapping the Formulation to the Research Questions}
\label{subsec:rq_mapping}

The formulation isolates the two contributions so that each research question
has a direct experimental test. Table~\ref{tab:rq_mapping} summarizes this
mapping.

\begin{table}[t]
\centering
\caption{Mapping between research questions, controlled comparisons, and
principal evidence.}
\label{tab:rq_mapping}
\small
\setlength{\tabcolsep}{5pt}
\renewcommand{\arraystretch}{1.15}
\begin{tabular}{@{}p{1.1cm}p{5.0cm}p{5.2cm}p{4.2cm}@{}}
\toprule
\textbf{RQ} & \textbf{Question in the formulation} & \textbf{Primary controlled comparisons} & \textbf{Principal evidence} \\
\midrule
RQ1 & Does $\mathbf{z}^{(1)}_i$ contain more useful attack information than aggregate flow statistics alone? & Flow-statistics MLP; CNN/LSTM/GRU baselines; positional/time encoding ablations; sequence-length and packet-selection studies; flow-fusion ablation. & Macro-F1, macro-precision, macro-recall, attack-class F1, PR-AUC, ROC-AUC, confusion matrix, and FPR. \\
RQ2 & Does $\mathcal{C}^{(r)}_i$ improve prediction beyond an independent classifier on $\mathbf{z}^{(1)}_i$? & Stage1-only and no-context controls; time-only, source-host, destination-host, and endpoint context; window-size ablation; Transformer, LSTM, GRU, and CNN--LSTM Stage2 encoders. & The same detection metrics, together with context availability and changes in FP and FN. \\
Cost & What computational cost is introduced by each representation and context model? & Stage1 and Stage2 models under the same implementation and hardware protocol. & Parameter count, training time, test inference time, and throughput where available. \\
\bottomrule
\end{tabular}
\end{table}

This separation is central to interpreting the final results. An improvement in
Stage1 supports the claim that packet timing and flow-statistical fusion improve
the representation of an individual flow. A further improvement from Stage2,
under the same Stage1 embeddings and test flows, supports the distinct claim
that related historical flows provide additional evidence. Efficiency results
then determine whether either improvement is operationally proportionate to its
computational cost.

\subsection{Chapter Summary}
\label{subsec:problem_summary}

The proposed problem is a causal, hierarchical, payload-independent
classification task. Stage1 maps an ordered and timed packet sequence together
with its flow statistics to a fused embedding. Stage2 constructs a bounded
historical context over these embeddings and predicts the label of the final
target flow. The separation between the stages supports controlled answers to
RQ1 and RQ2, while validation-only threshold selection, mask-aware sequence
processing, and explicit context policies establish the constraints required
for a reliable evaluation. The following chapters describe how the PCAP data
are transformed into the variables defined here and how the two model stages
are implemented and evaluated.
